\documentclass[11pt]{article}
\usepackage[margin=1.1in]{geometry}
\usepackage{amsmath,amssymb,amsthm}
\usepackage[colorlinks=true,linkcolor=blue,citecolor=blue]{hyperref}

\newtheorem{theorem}{Theorem}[section]
\newtheorem{lemma}[theorem]{Lemma}
\newtheorem{corollary}[theorem]{Corollary}
\newtheorem{remark}[theorem]{Remark}
\newtheorem{conjecture}[theorem]{Conjecture}

\newcommand{\R}{\mathbb{R}}
\newcommand{\E}{\mathbb{E}}
\newcommand{\PP}{\mathbb{P}}
\newcommand{\emp}{\bar\mu_N}
\newcommand{\vol}{\operatorname{vol}}

\title{The edge exponent of the uniform doubling condition\\
and an interpolation conjecture}
\author{Yitzchak Shmalo\\
\small Einstein Institute of Mathematics, The Hebrew University of Jerusalem, Jerusalem, Israel\\
\small }
\date{June 2026}

\begin{document}
\maketitle

\begin{abstract}
In a sequence of companion papers we determined the critical scaling of the
uniform doubling condition of
[Shmalo, \emph{Ann.\ Math.\ Artif.\ Intell.} (2024)] for i.i.d.\ samples in
two regimes: $\Theta((\log N/N)^{2/(n+1)})$ for pure slabs over smooth convex
bodies, and $\Theta((\log N/N)^{1/n})$ in the presence of conical truncations.
This note locates the first intermediate exponent.  For the unit ball
truncated by a \emph{single transversal half-space} --- so that the boundary
of the truncation region has an $(n-2)$-dimensional edge but no
lower-dimensional strata --- we prove that with probability $1-o(1)$ the
doubling inequality fails at every grid scale below
$c\,(\log N/N)^{2/(n+2)}$, $n\ge3$.  The obstruction is the most isolated
sample point of the \emph{edge caps}, whose volumes scale as
$h^{(n+2)/2}$, intermediate between the smooth cap exponent
$\frac{n+1}2$ and the conical exponent $n$.  A second-moment argument over
$\asymp(N/\log N)^{(n-2)/(n+2)}$ congruent caps along the edge upgrades the
failure probability to $1-o(1)$, with the logarithm.  We conjecture the general
interpolation law: for truncation families whose regions have all boundary
strata of dimension at least $j$ and transversal intersections, the critical
exponent is $\frac{2}{2n-j}$ --- recovering
$\frac{2}{n+1}$ for smooth boundaries ($j=n-1$), $\frac{2}{n+2}$ for edges
($j=n-2$), and $\frac1n$ for conical points ($j=0$) --- and we explain what a
matching upper bound requires.
\end{abstract}

\section{Setting and result}\label{sec:intro}

Let $K=B^n\subset\R^n$ be the closed unit ball, $n\ge3$, $\mu$ the uniform
probability measure on $K$, and $T=\{X_1,\dots,X_N\}$ i.i.d.\ samples with
empirical measure $\emp$.  Truncated slabs
$S_w=\{x:|\langle x-s,u\rangle|\le\frac w2\}\cap H$ and the uniform doubling
condition (UDC) with constants $(\kappa,\delta,\sigma,\beta,\ell)$ are as in
\cite{ShmaloUDC,companion2}.  The critical scaling of the smallest admissible
scaling factor $\ell$ is known in two regimes
\cite{companion3,companion4,companion5}: exactly
$\Theta((\log N/N)^{2/(n+1)})$ for pure slabs ($k=0$) over bodies with a
rolling ball, and exactly $\Theta((\log N/N)^{1/n})$ for thick conical
truncation geometry.  Here we consider the geometrically intermediate case:
a single half-space truncation
\[
H:=\{x:\ x_1\le0\},
\qquad
G:=K\cap H,
\]
whose boundary consists of two smooth $(n-1)$-dimensional pieces (a spherical
cap and a flat disk) meeting transversally (at dihedral angle $\pi/2$) along
the $(n-2)$-dimensional edge
\[
E\;:=\;\partial K\cap\partial H\;=\;\{x:\ |x|=1,\ x_1=0\}.
\]
The region $G$ is uniformly thick (it contains the ball
$B(-\tfrac12e_1,\tfrac14)$), and the upper bound of \cite{companion2} applies:
the UDC holds with probability $1-\frac1N$ at scaling
$\ell_N=O((\log N/N)^{1/n})$.  We prove that the true critical scale for this
single truncation is governed by a new exponent.

\begin{theorem}\label{thm:main}
Let $n\ge3$ and fix $\kappa>1$, $\sigma>0$, $\beta>0$, $\delta\ge\frac3N$.
There are constants $\varepsilon_2(n,\kappa)>0$ and $c(n,\kappa)>0$ such that,
with probability $1-o(1)$ as $N\to\infty$, the doubling inequality fails for
some slab truncated by the single truncation set $\{(e_1,0)\}$ at every grid
scale below $c\,(\log N/N)^{2/(n+2)}$.  Consequently, with probability
$1-o(1)$, the UDC for this truncation requires
\[
\ell\;\ge\;c\,\Bigl(\frac{\log N}{N}\Bigr)^{2/(n+2)} .
\]
\end{theorem}

Since $\frac{2}{n+2}>\frac1n$ would fail and indeed
$\frac{2}{n+2}<\frac{2}{n+1}$, the bound interpolates strictly between the two
known regimes: a single transversal truncation of a smooth body already
destroys the smooth exponent $\frac{2}{n+1}$, but does not produce the conical
exponent $\frac1n$ --- for that, by \cite{companion3}, genuinely conical
(lowest-dimensional) boundary strata are needed.  We record the natural guess
suggested by the three data points.

\begin{conjecture}[Interpolation law]\label{conj:law}
For i.i.d.\ samples with density bounded above and below on a convex body, and
truncation families of bounded size whose nonempty regions are uniformly thick
and have all boundary strata transversal and of dimension at least $j$, the
critical scaling of the uniform doubling condition is
\[
\ell_N\;\asymp\;\Bigl(\frac{\log N}{N}\Bigr)^{2/(2n-j)} ,
\]
where for $j=n-1$ (smooth, rolling-ball boundary) this is the proven
$\frac{2}{n+1}$ \cite{companion4}, for $j=n-2$ (transversal edges) it is the
exponent $\frac{2}{n+2}$ of Theorem~\ref{thm:main}, and for $j=0$ (conical
points) it is the proven $\frac1n$ \cite{companion3,companion5}.
\end{conjecture}

The heuristic behind the law is uniform: the critical resolution is the
isolation depth of the most isolated extremal sample point, and a cap of depth
$h$ at a $j$-dimensional transversal stratum has volume
$\asymp h^{\,(2n-j)/2}$ --- each of the $n-1-j$ transverse flat directions
contributes a factor $h$, the normal direction contributes $h$, and each of
the $j$ curved stratum directions contributes $\sqrt h$.  Theorem
\ref{thm:main} verifies the count for $j=n-2$:
$h^{2}\cdot h^{(n-2)/2}=h^{(n+2)/2}$.

\section{Edge caps}\label{sec:caps}

Fix the angle $\varphi=\frac\pi4$ and, for a point $p\in E$, the direction
\[
u_p\;:=\;\tfrac1{\sqrt2}\,e_1+\tfrac1{\sqrt2}\,p .
\]
Note $\langle x,u_p\rangle$ is maximized over $G$ at $x=p$, with maximum
$\tfrac1{\sqrt2}$: for $x\in G$,
$\langle x,u_p\rangle=\tfrac1{\sqrt2}\bigl(x_1+\langle x,p\rangle\bigr)
\le\tfrac1{\sqrt2}\bigl(0+|x|\,|p|\bigr)\le\tfrac1{\sqrt2}$,
with equality iff $x_1=0$ and $x=p$.  Define the \emph{edge cap}
\[
C_p(h)\;:=\;\Bigl\{x\in G:\ \langle x,u_p\rangle\ge\tfrac1{\sqrt2}-h\Bigr\} .
\]

\begin{lemma}[Edge cap volume]\label{lem:cap}
There are constants $0<c_8(n)\le c_9(n)$ such that for all $p\in E$ and
$0<h\le h_0(n)$,
\[
c_8\,h^{\frac{n+2}2}\;\le\;\mu\bigl(C_p(h)\bigr)\;\le\;c_9\,h^{\frac{n+2}2}.
\]
\end{lemma}

\begin{proof}
By rotational symmetry about the $e_1$-axis we may take $p=e_2$.  Write
$x=(x_1,x_2,x'')$ with $x''\in\R^{n-2}$, and $h':=\sqrt2\,h$.  The cap is
\[
C_p(h)=\Bigl\{x:\ |x|\le1,\ x_1\le0,\ x_2-|x_1|\ge1-h'\Bigr\}
\]
(using $x_1\le0$ to write $x_1=-|x_1|$).  On $|x|\le1$ we have
$x_2\le\sqrt{1-x_1^2-|x''|^2}\le1-\tfrac{x_1^2+|x''|^2}2$, and
$x_2\ge1-h'+|x_1|$ forces
\begin{equation}\label{eq:constraints}
|x_1|\;\le\;h',\qquad |x''|^2\;\le\;2h',\qquad
1-h'+|x_1|\;\le\;x_2\;\le\;\sqrt{1-x_1^2-|x''|^2}.
\end{equation}
\emph{Upper bound:} the $x_2$-fibre has length at most
$\bigl(1-\tfrac{x_1^2+|x''|^2}2\bigr)-(1-h'+|x_1|)\le h'$, so
$\vol(C_p(h))\le h'\cdot(2h')\cdot\omega_{n-2}(2h')^{\frac{n-2}2}
\le C(n)h^{\frac{n+2}2}$, where the three factors bound the
$x_2$-, $x_1$- and $x''$-ranges.  \emph{Lower bound:} restrict to
$|x_1|\le\frac{h'}{8}$, $|x''|^2\le\frac{h'}2$; then, for
$h'\le h_0$ small, using $\sqrt{1-t}\ge1-t$ for $t\in[0,1]$,
\[
\sqrt{1-x_1^2-|x''|^2}-(1-h'+|x_1|)
\;\ge\;h'-|x_1|-x_1^2-|x''|^2
\;\ge\;h'-\tfrac{h'}8-\tfrac{h'}{64}-\tfrac{h'}2\;\ge\;\tfrac{h'}4,
\]
so the fibre length is at least $\frac{h'}4$ on a set of
$(x_1,x'')$-measure $\frac{h'}4\cdot\omega_{n-2}(\frac{h'}2)^{\frac{n-2}2}$,
giving $\vol(C_p(h))\ge c(n)h^{\frac{n+2}2}$.  Dividing by
$\vol(B^n)$ gives the bounds for $\mu$.
\end{proof}

\begin{lemma}[Footprint and disjointness]\label{lem:disjoint}
There is $c_{10}(n)$ such that
$C_p(h)\subseteq\{x:\ |x-p|\le c_{10}\sqrt h\}$ for $h\le h_0$.  Consequently,
if $p_1,\dots,p_M\in E$ are pairwise at distance greater than
$2c_{10}\sqrt H$, the caps $C_{p_j}(H)$ are pairwise disjoint, and on the
$(n-2)$-dimensional sphere $E$ one may choose
\begin{equation}\label{eq:M}
M\;\ge\;c_{11}(n)\,H^{-\frac{n-2}2}
\end{equation}
such points.
\end{lemma}

\begin{proof}
From \eqref{eq:constraints} (with $p=e_2$): $|x_1|\le h'$,
$|x''|\le\sqrt{2h'}$ and $1\ge x_2\ge1-h'+|x_1|\ge1-h'$, so
$|x-e_2|^2\le h'^2+2h'+h'^2\le4h'$.  The packing bound on the
$(n-2)$-dimensional sphere $E$ --- of dimension $n-2\ge1$ for $n\ge3$ --- is
the standard volumetric estimate.
\end{proof}

\begin{lemma}[Stall]\label{lem:stall}
Fix $\kappa>1$, $\sigma>0$, $\delta\ge\frac3N$, $p\in E$, and depths
$0<h\le\frac H4$, $H\le h_0$.  Suppose exactly one sample point lies in
$C_p(H)$ and that point, $X^*$, lies in $C_p(h)$.  Then for every
$w\le H/\kappa$ the slab with direction $u_p$, center $X^*$, and truncation
$\{(e_1,0)\}$ satisfies
$\emp(S_w)=\emp(S_{\kappa w})=\frac1N<\min\{\delta,(1+\sigma)\emp(S_w)\}$;
hence the doubling inequality fails at every scale $w\le H/\kappa$, and the
UDC with scaling factor $\ell\le H/\kappa^2$ fails.
\end{lemma}

\begin{proof}
For $w'\le H$:
$S_{w'}\subseteq G\cap\{\langle x,u_p\rangle\ge\langle X^*,u_p\rangle
-\frac{w'}2\}\subseteq C_p\bigl(h+\frac{w'}2\bigr)\subseteq C_p(H)$, since
$h+\frac{w'}2\le\frac H4+\frac H2<H$; the rest is as in
\cite{companion4,companion5}.
\end{proof}

\section{Proof of Theorem \ref{thm:main}}\label{sec:proof}

\begin{proof}
Set $q:=(4\kappa^2)^{\frac{n+2}2}\frac{c_9}{c_8}$,
$\varepsilon_2:=\frac{n-2}{2(n+2)\,q}$ and $\tau_N:=\varepsilon_2\log N$.
Define $h_N$ by $\mu(C_p(h_N))=\frac{\tau_N}N$ (the cap mass is continuous and
strictly increasing in $h$, and independent of $p$ by symmetry); by
Lemma~\ref{lem:cap}, $h_N\asymp(\tau_N/N)^{2/(n+2)}$.  Set
$H_N:=4\kappa^2h_N$, so by Lemma~\ref{lem:cap},
$\mu(C_p(H_N))\le q\,\frac{\tau_N}N$.

Choose $p_1,\dots,p_M\in E$ as in Lemma~\ref{lem:disjoint} with
$M\ge c_{11}H_N^{-\frac{n-2}2}\ge c_{12}(N/\tau_N)^{\frac{n-2}{n+2}}$, and let
\[
E_j:=\Bigl\{\#\bigl(T\cap C_{p_j}(H_N)\bigr)=1,\
T\cap C_{p_j}(H_N)\subseteq C_{p_j}(h_N)\Bigr\},
\qquad
X:=\sum_{j\le M}\mathbf1_{E_j}.
\]
Exactly as in \cite{companion4}:
\[
p:=\PP(E_j)=N\cdot\frac{\tau_N}N\bigl(1-\mu(C_{p_j}(H_N))\bigr)^{N-1}
\ge\tau_N\,e^{-2q\tau_N}=\tau_N\,N^{-2q\varepsilon_2}
=\tau_N\,N^{-\frac{n-2}{n+2}},
\]
so $\E X=Mp\ge c_{12}\,\tau_N^{1-\frac{n-2}{n+2}}\to\infty$; replacing
$\varepsilon_2$ by $\varepsilon_2/2$ gives the safer
$\E X\ge c_{12}N^{\frac{n-2}{2(n+2)}}\to\infty$.  The caps being disjoint, the
two-cap correlation ratio is
$\exp(O(N\mu(C(H_N))^2))=\exp(O((\log N)^2/N))=1+o(1)$ as in
\cite{companion4}, so $\PP(X=0)\le\frac1{\E X}+o(1)=o(1)$ by Chebyshev.
With probability $1-o(1)$ some $E_j$ occurs, and Lemma~\ref{lem:stall} forces
$\ell>H_N/\kappa^2=4h_N\ge c(n,\kappa)(\log N/N)^{2/(n+2)}$.
\end{proof}

\section{Towards the matching upper bound}\label{sec:upper}

Conjecture~\ref{conj:law} predicts that for the single transversal truncation
of Theorem~\ref{thm:main} the UDC in fact \emph{holds} with high probability at
scaling $\ell_N=O((\log N/N)^{2/(n+2)})$.  By the proof skeleton of
\cite{companion2,companion4}, this reduces to a deterministic geometric
estimate: a uniform lower bound
\[
\vol\Bigl(G\cap\bigl\{y:\ |\langle y-x,u\rangle|\le L\bigr\}\Bigr)
\;\ge\;c\,L^{\frac{n+2}2}
\qquad(x\in G,\ u\in S^{n-1},\ 0<L\le L_0),
\]
the analogue for $G=K\cap H$ of the rolling-ball slab-volume bound of
\cite{companion4}.  The anchor--unimodality framework of \cite{companion4}
reduces this to lower bounds on the level-section volumes
$A(\tau)$ of $G$ at depth $h$ below the $u$-extreme, for every direction: the
required anchor value is $A\gtrsim h^{n/2}$ at some depth $h\le L/4$.  Three
regimes occur: extremes on the smooth spherical part (where the rolling-ball
sections give the stronger rate $h^{(n-1)/2}$, but the half-space may cut
them), extremes on the flat disk (where the sections contain
$(n-1)$-dimensional disk caps of volume $\asymp h^{n/2}$), and extremes on the
edge (where the explicit computation of Lemma~\ref{lem:cap} gives sections of
volume $\asymp h^{n/2}$).  A convex-segment witness from the extreme point to
the deepest point of the cut rolling ball shows that the section's penetration
depth into $H$ grows at least linearly in $h$, which controls the cut in the
first regime.  Assembling these into a single uniform constant --- in
particular through the transition directions where the extreme point migrates
from the spherical part across the edge to the flat disk --- requires a
careful case analysis that we have not completed, and which we leave, together
with the general law of Conjecture~\ref{conj:law}, as the open problem of this
note.  We emphasize that all three regime computations individually produce
the predicted exponent $\frac{n+2}2=\frac{2n-j}2$ with $j=n-2$, in agreement
with the stratum-counting heuristic.

\subsection*{Acknowledgments}
The research presented in this paper was supported by the European Research
Council (ERC) under the European Union's Horizon Europe research and
innovation programme (grant agreement No.\ 101041711), by the Israel Science
Foundation (grant number 2258/19) and by the Israel Science Foundation (ISF
Grant 4101/25).

\subsection*{Declaration on generative AI}
During the preparation of this work the author used generative AI tools to
accelerate drafting and revision.  All mathematical claims, proofs, and
bibliographic information were subsequently reviewed and edited by the author,
who takes full responsibility for the content of the manuscript.

\begin{thebibliography}{9}

\bibitem{ShmaloUDC}
Y.~Shmalo,
\emph{Stability of accuracy for the training of DNNs via the uniform doubling
condition},
Ann.\ Math.\ Artif.\ Intell.\ (2024), DOI 10.1007/s10472-023-09919-1.

\bibitem{companion2}
Y.~Shmalo,
\emph{The uniform doubling condition holds with high probability for
independent samples}, manuscript, 2026.

\bibitem{companion3}
Y.~Shmalo,
\emph{The critical scaling of the uniform doubling condition: smooth bodies
versus truncations}, manuscript, 2026.

\bibitem{companion4}
Y.~Shmalo,
\emph{The logarithm is necessary: exact critical scaling of the uniform
doubling condition on the ball}, manuscript, 2026.

\bibitem{companion5}
Y.~Shmalo,
\emph{The exact critical scaling of the uniform doubling condition with
truncations}, manuscript, 2026.

\end{thebibliography}

\end{document}
